\begin{explanationbox}
	\textbf{\textcolor{CExplanation}{一句话直觉}}

	二次函数在固定区间上的最值，取决于对称轴相对于区间的位置以及开口方向；位置不同取法不同，只需分类即可。

	\medskip
	\textbf{\textcolor{CExplanation}{核心拆解}}
	\begin{description}[style=nextline, leftmargin=2.8em, labelsep=0.5em, itemsep=0.45em]
		\item[\textbf{要点一（开口定向）：}] 由 $a$ 的正负决定抛物线开口方向，$a>0$ 开口向上，$a<0$ 开口向下，单调性相反。
		\item[\textbf{要点二（对称轴分段）：}] 最小值（$a>0$）分三段：对称轴在区间左、内、右，对应取左端点、顶点、右端点。
		\item[\textbf{要点三（中点判远）：}] 最大值（$a>0$）仅分两段：对称轴在区间中点左侧时右端点更远（值更大），在中点右侧时左端点更远。
	\end{description}

	\medskip
	\textbf{\textcolor{CExplanation}{几何本质（对称轴与区间）}}

	开口向上时，离对称轴越远的点函数值越大。最小值是区间内距对称轴最近的点（可能端点或顶点），最大值是距对称轴最远的端点。开口向下则相反。

	\medskip
	\textbf{\textcolor{CExplanation}{代数意义（分类讨论）}}

	将区间端点与对称轴的具体数量关系转化为函数值的大小比较，通过解不等式或等式确定最值表达式的范围。

	\medskip
	\textbf{\textcolor{CExplanation}{考点价值（分类与对称）}}
	\begin{itemize}[leftmargin=*, itemsep=0.5em, topsep=0.4em]
		\item \textbf{考法一（直接套用）：} 给定具体二次函数与区间，直接按 $a>0$ 或 $a<0$ 套用公式求出最值。
		\item \textbf{考法二（含参讨论）：} 二次函数含参数，需按对称轴与区间关系分类，写出最值关于参数的分段函数。
		\item \textbf{考法三（反向求参）：} 已知最值条件，反求参数值或范围，需分类解方程并检验。
	\end{itemize}

	\medskip
	\textbf{\textcolor{CExplanation}{顿悟点}}

	所谓“定区间动轴”，分类标准就是对称轴是否落在区间内，以及对称轴与区间中点的左右关系；这是所有此类问题共同的模型。

	\medskip
	\textbf{\textcolor{CExplanation}{使用场景}}
	\begin{itemize}[leftmargin=*, itemsep=0.5em, topsep=0.4em]
		\item \textbf{场景一（定区间动轴）：} 二次项系数为参数或一次项系数为参数，对称轴随参数变化，区间固定。
		\item \textbf{场景二（已知最值求参）：} 给出函数在区间上的最值，求参数的值或范围，需要分类讨论后解方程。
	\end{itemize}
\end{explanationbox}
